\chapter{Oppervlaktes en volumes}\label{ch:g7:areas}

In jaar 6 maten we rechthoeken, \omterm{def:g6:shapes:triangles}{rechthoekige driehoeken} en \omterm{def:g5:solids:def}{balken}
(\cref{ch:g6:measure}). Met een schaar in de hand --- hier een stuk afknippen,
daar weer opplakken --- geven dezelfde ideeën nu de \omterm{def:g6:measure:area}{oppervlaktes} van
parallellogrammen, van \emph{alle} driehoeken en van de schijf, en de \omterm{def:g6:measure:volume}{volumes} van
\omterm{def:g7:areas:prism}{prisma}'s en \omterm{def:g7:areas:prism}{cilinders}.

\section{Oppervlakte van een parallellogram}

\begin{theorem}[Parallellogram]\label{thm:g7:areas:parallelogram}
Een \omterm{def:g7:central:parallelogram}{parallellogram} met een zijde van lengte $b$ (de \emph{basis}) en afstand $h$
tussen die zijde en de zijde ertegenover (de \emph{hoogte}) heeft \omterm{def:g6:measure:area}{oppervlakte}
\[
A = b \times h .
\]
\end{theorem}

\begin{proof}[Bewijs met de schaar]
Knip aan de ene kant de \omterm{def:g6:shapes:triangles}{rechthoekige driehoek} af die uitsteekt en plak hem aan de
andere kant: het \omterm{def:g7:central:parallelogram}{parallellogram} wordt een rechthoek $b \times h$ met dezelfde
\omterm{def:g6:measure:area}{oppervlakte}.
\end{proof}

\begin{omfigure}
\begin{tikzpicture}[scale=0.82]
  \fill[omDef!14] (0,0) -- (4,0) -- (5.2,1.8) -- (1.2,1.8) -- cycle;
  \draw[thick] (0,0) -- (4,0) -- (5.2,1.8) -- (1.2,1.8) -- cycle;
  \fill[omThm!25] (0,0) -- (1.2,1.8) -- (1.2,0) -- cycle;
  \draw[omThm, thick] (0,0) -- (1.2,1.8) -- (1.2,0) -- cycle;
  \draw[dashed] (1.2,0) -- (1.2,1.8);
  \node[below, font=\small] at (2.0,0) {$b$};
  \draw[Stealth-Stealth] (6.0,0) -- (6.0,1.8)
    node[midway, right, font=\small] {$h$};
  \begin{scope}[xshift=8.2cm]
  \fill[omDef!14] (0,0) rectangle (4,1.8);
  \draw[thick] (0,0) rectangle (4,1.8);
  \fill[omThm!25] (2.8,0) -- (4,1.8) -- (4,0) -- cycle;
  \draw[omThm, thick] (2.8,0) -- (4,1.8) -- (4,0) -- cycle;
  \node[below, font=\small] at (2.0,0) {$b$};
  \end{scope}
\end{tikzpicture}

{\small Het bewijs met de schaar: de rode driehoek van het linkeruiteinde naar het
rechteruiteinde schuiven maakt van het \omterm{def:g7:central:parallelogram}{parallellogram} een rechthoek --- dezelfde
\omterm{def:g6:measure:area}{oppervlakte}, $b \times h$.}
\end{omfigure}

\begin{remark}
De hoogte is de \emph{loodrechte} afstand tussen de twee bases, en niet de lengte
van de schuine zijde! Een sterk scheef \omterm{def:g7:central:parallelogram}{parallellogram} kan lange zijden en een kleine
\omterm{def:g6:measure:area}{oppervlakte} hebben.
\end{remark}

\section{Oppervlakte van een driehoek}

\begin{theorem}[Driehoek]\label{thm:g7:areas:triangle}
Een driehoek met basis $b$ en bijbehorende hoogte $h$ (de loodrechte afstand van het
\omterm{def:g5:solids:def}{hoekpunt} ertegenover tot de basislijn) heeft \omterm{def:g6:measure:area}{oppervlakte}
\[
A = \frac{b \times h}{2} .
\]
\end{theorem}

\begin{proof}
Twee kopieën van de driehoek, waarvan er één een halve draai om het midden van een
zijde gedraaid is, passen samen in een \omterm{def:g7:central:parallelogram}{parallellogram} met basis $b$ en hoogte $h$
(dat is de \omterm{def:g7:central:def}{puntspiegeling} aan het werk, \cref{prop:g7:central:preserved}). De
driehoek is de \omterm{def:g3:division:half}{helft} daarvan: $\frac{bh}{2}$.
\end{proof}

\begin{omfigure}
\begin{tikzpicture}[scale=0.9]
  \fill[omDef!14] (0,0) -- (4,0) -- (1.4,2.1) -- cycle;
  \draw[thick] (0,0) -- (4,0) -- (1.4,2.1) -- cycle;
  \fill[omThm!12] (4,0) -- (1.4,2.1) -- (5.4,2.1) -- cycle;
  \draw[thick, omThm, dashed] (4,0) -- (5.4,2.1) -- (1.4,2.1);
  \draw[dashed] (1.4,0) -- (1.4,2.1);
  \draw[thin] (1.4,0) ++(-0.2,0) -- ++(0,0.2) -- ++(0.2,0);
  \node[below, font=\small] at (2,0) {$b$};
  \node[right, font=\small] at (1.42,1.0) {$h$};
\end{tikzpicture}

{\small Een driehoek met zijn halfgedraaide kopie (gestippeld) vormt een
\omterm{def:g7:central:parallelogram}{parallellogram}: de \omterm{def:g6:measure:area}{oppervlakte} van de driehoek is $\frac{b\times h}{2}$. Elk van de
drie zijden kan als basis dienen --- met \emph{haar eigen} hoogte.}
\end{omfigure}

\begin{example}\label{ex:g7:areas:triangle}
Een driehoek heeft een basis van $9$ cm en de bijbehorende hoogte meet $4$ cm:
\[
A = \frac{9 \times 4}{2} = \frac{36}{2} = 18 \text{ cm}^2 .
\]
Bij een \emph{rechthoekige} driehoek zijn de twee rechthoekszijden een basis met haar
hoogte: zo vinden we de formule van \cref{prop:g6:measure:formulas} terug.
\end{example}

\section{Oppervlakte van een schijf}

\begin{theorem}[Schijf]\label{thm:g7:areas:disk}
Een schijf met \omterm{def:g3:shapes:shapes}{radius} $r$ heeft \omterm{def:g6:measure:area}{oppervlakte}
\[
A = \pi r^2 .
\]
\end{theorem}

\begin{proof}[Idee van het bewijs]
Knip de schijf in heel veel dunne gelijke partjes en leg ze kop aan staart: ze vormen
bijna een \omterm{def:g7:central:parallelogram}{parallellogram} met hoogte $r$ waarvan de basis de \omterm{def:g3:division:half}{helft} van de \omterm{def:g6:measure:perimeter}{omtrek} is,
$\pi r$ (\cref{prop:g6:measure:circle}). Zijn \omterm{def:g6:measure:area}{oppervlakte} ligt dicht bij
$\pi r \times r = \pi r^2$, en de benadering wordt perfect als de partjes dunner
worden. Een volledig bewijs heeft integraalrekening nodig, uit het bovenbouwvolume.
\end{proof}

\begin{omfigure}
\begin{tikzpicture}[scale=0.85]
  \foreach \a in {0,45,...,315}
    \fill[omDef!20, draw=omDef] (0,0) -- (\a:1.4)
      arc (\a:\a+45:1.4) -- cycle;
  \begin{scope}[xshift=4.3cm, yshift=-0.15cm]
  \foreach \i in {0,1,2,3} {
    \fill[omDef!20, draw=omDef]
      ({\i*1.1},0) -- ({\i*1.1+0.55},1.4) -- ({\i*1.1+1.1},0) -- cycle;
    \fill[omDef!20, draw=omDef]
      ({\i*1.1+0.55},1.4) -- ({\i*1.1+1.1},0) -- ({\i*1.1+1.65},1.4)
      -- cycle;
  }
  \draw[Stealth-Stealth] (5.2,0) -- (5.2,1.4)
    node[midway, right, font=\small] {$\approx r$};
  \node[below, font=\small] at (2.5,-0.1) {basis $\approx \pi r$};
  \end{scope}
\end{tikzpicture}

{\small De schijf in partjes knippen en die uitrollen: bijna een \omterm{def:g7:central:parallelogram}{parallellogram} met
basis $\pi r$ (de \omterm{def:g3:division:half}{helft} van de rand) en hoogte $r$.}
\end{omfigure}

\begin{example}\label{ex:g7:areas:disk}
Een rond tafelblad heeft \omterm{def:g3:shapes:shapes}{radius} $60$ cm. Zijn \omterm{def:g6:measure:area}{oppervlakte} is
$\pi \times 60^2 = 3600\pi \approx 11\,310$ cm$^2$, iets meer dan $1.1$ m$^2$. Let
op het \omterm{ex:g1:subtraction:difference}{verschil}: de \emph{\omterm{def:g6:measure:perimeter}{omtrek}} $2\pi r$ gebruikt $r$, de \emph{\omterm{def:g6:measure:area}{oppervlakte}}
$\pi r^2$ gebruikt $r^2$.
\end{example}

\section{Prisma's en cilinders}

\begin{definition}[Prisma, cilinder]\label{def:g7:areas:prism}
Een \emph{prisma}\index{prisma} is een \omterm{def:g5:solids:def}{lichaam} met twee identieke parallelle
veelhoekige \omterm{def:g5:solids:def}{zijvlakken} (de \emph{grondvlakken}), verbonden door rechthoeken; een
\emph{cilinder}\index{cilinder} is hetzelfde met schijven als grondvlakken. De
\emph{hoogte} is de afstand tussen de twee grondvlakken.
\end{definition}

\begin{theorem}[Volumes]\label{thm:g7:areas:volume}
Voor een \omterm{def:g7:areas:prism}{prisma} of een \omterm{def:g7:areas:prism}{cilinder} met grondvlakoppervlakte $B$ en hoogte $h$:
\[
V = B \times h .
\]
\end{theorem}
\admitted

\begin{example}\label{ex:g7:areas:volume}
Een tent is een \omterm{def:g7:areas:prism}{prisma} dat op zijn zij ligt: haar grondvlakken zijn driehoeken met
basis $2$ m en hoogte $1.5$ m, en de tent is $3$ m lang.
\begin{enumerate}
  \item \omterm{def:g6:measure:area}{Oppervlakte} van het grondvlak: $B = \frac{2 \times 1.5}{2} = 1.5$ m$^2$.
  \item \omterm{def:g6:measure:volume}{Volume}: $V = B \times h = 1.5 \times 3 = 4.5$ m$^3$.
\end{enumerate}
Een conservenblik met \omterm{def:g3:shapes:shapes}{radius} $4$ cm en hoogte $11$ cm:
$V = \pi \times 4^2 \times 11 = 176\pi \approx 553$ cm$^3$, ruwweg een halve liter.
\end{example}

\section{Oefeningen}

\begin{exercise}[$\star$]\label{exo:g7:areas:1}
Reken de \omterm{def:g6:measure:area}{oppervlakte} uit van een \omterm{def:g7:central:parallelogram}{parallellogram} met basis $8$ cm en hoogte $5$ cm;
en van een ander met basis $6.5$ m en hoogte $4$ m.
\end{exercise}

\begin{exercise}[$\star$]\label{exo:g7:areas:2}
Een \omterm{def:g7:central:parallelogram}{parallellogram} heeft zijden van $10$ cm en $6$ cm, en de hoogte bij de basis van
$10$ cm meet $4$ cm. Reken zijn \omterm{def:g6:measure:area}{oppervlakte} uit. Waarom is het antwoord geen $60$
cm$^2$?
\end{exercise}

\begin{exercise}[$\star$]\label{exo:g7:areas:3}
Reken de \omterm{def:g6:measure:area}{oppervlakte} uit van een driehoek met basis $12$ cm en hoogte $7$ cm; van een
\omterm{def:g6:shapes:triangles}{rechthoekige driehoek} met rechthoekszijden $9$ en $6$; en van een driehoek met basis
$5.5$ m en hoogte $2$ m.
\end{exercise}

\begin{exercise}[$\star$]\label{exo:g7:areas:4}
Reken de \omterm{def:g6:measure:area}{oppervlakte} uit van een schijf met \omterm{def:g3:shapes:shapes}{radius} $5$ cm, en daarna van een halve
schijf met \omterm{def:g3:shapes:shapes}{radius} $10$ cm (precieze waarden met $\pi$, en daarna afgerond op de
eenheid).
\end{exercise}

\begin{exercise}[$\star$]\label{exo:g7:areas:5}
Reken de \omterm{def:g6:measure:perimeter}{omtrek} \emph{en} de \omterm{def:g6:measure:area}{oppervlakte} uit van een schijf met \omterm{def:g3:shapes:shapes}{radius} $3$ cm. Welk
van de twee antwoorden is in cm en welk in cm$^2$?
\end{exercise}

\begin{exercise}[$\star$]\label{exo:g7:areas:6}
Reken het \omterm{def:g6:measure:volume}{volume} uit van een \omterm{def:g7:areas:prism}{prisma} met grondvlakoppervlakte $24$ cm$^2$ en hoogte
$10$ cm; en van een \omterm{def:g7:areas:prism}{cilinder} met \omterm{def:g3:shapes:shapes}{radius} $3$ cm en hoogte $8$ cm (precies, en daarna
afgerond).
\end{exercise}

\begin{exercise}[$\star$]\label{exo:g7:areas:7}
Een driehoek heeft \omterm{def:g6:measure:area}{oppervlakte} $28$ cm$^2$ en basis $8$ cm. Zoek de bijbehorende
hoogte, en schrijf eerst de vergelijking op.
\end{exercise}

\begin{exercise}[$\star\star$]\label{exo:g7:areas:8}
Teken een willekeurige driehoek en meet zorgvuldig de drie paren basis en hoogte.
Reken voor elk paar $\frac{b \times h}{2}$ uit: de drie uitkomsten moeten
overeenkomen (op de meetfout na). Waarom?
\end{exercise}

\begin{exercise}[$\star\star$]\label{exo:g7:areas:9}
Een tuin in de vorm van een trapezium kun je met één diagonaal in twee driehoeken
verdelen met dezelfde hoogte $h = 20$ m en bases $30$ m en $18$ m. Reken haar totale
\omterm{def:g6:measure:area}{oppervlakte} uit. Kun je een algemene formule voor het trapezium raden?
\end{exercise}

\begin{exercise}[$\star\star$]\label{exo:g7:areas:10}
Een cilindervormige vaas met \omterm{def:g3:shapes:shapes}{radius} $6$ cm bevat water tot een hoogte van $15$ cm. Al
het water wordt in een lege \omterm{def:g5:solids:def}{balk} met een grondvlak van $18$ cm $\times$ $12$ cm
gegoten. Welke waterhoogte bereikt het daar? (Hetzelfde \omterm{def:g6:measure:volume}{volume}, een ander grondvlak.)
\end{exercise}

\begin{exercise}[$\star\star$]\label{exo:g7:areas:11}
Een pizza met \omterm{def:g4:geometry:circle}{diameter} $30$ cm kost $9$ euro; een met \omterm{def:g4:geometry:circle}{diameter} $40$ cm kost $14$
euro. Reken de twee \omterm{def:g6:measure:area}{oppervlaktes} uit en de prijs per $100$ cm$^2$. Welke pizza is
voordeliger?
\end{exercise}

\begin{exercise}[$\star\star\star$]\label{exo:g7:areas:12}
De twee rechthoekszijden van een \omterm{def:g6:shapes:triangles}{rechthoekige driehoek} meten $6$ cm en $8$ cm, en
zijn schuine zijde meet $10$ cm. Reken zijn \omterm{def:g6:measure:area}{oppervlakte} uit met de rechthoekszijden,
en gebruik die \omterm{def:g6:measure:area}{oppervlakte} daarna om de hoogte bij de \emph{schuine zijde} te vinden.
\end{exercise}

\section{Opgave: het verdwenen vierkantje}

\begin{problem}[{Weekendopgave --- een beroemd ``bewijs'' dat $64 = 65$, ontmaskerd
met oppervlaktes en hellingen, met pizza's en ringen als nagerecht}]
\label{pb:g7:areas:1}
Een goochelaar knipt een vierkant van $8 \times 8$ in vier stukken, schuift ze wat
heen en weer, en legt ze weer samen tot een rechthoek van $5 \times 13$. Het publiek
rekent: $8 \times 8 = 64$, maar $5 \times 13 = 65$ --- er is uit het niets een
vierkantje verschenen! De oppervlakteformules van dit hoofdstuk zijn precies het
gereedschap dat de truc betrapt. Daarna zet de opgave eerlijk
oppervlakteredeneren aan het werk: op trapezia, op scheve stapels, op de economie van
pizza's --- en tot slot komt de goochelaar terug voor een tweede, nog vreemdere
voorstelling.

\medskip
\noindent\textbf{Deel I --- De truc, uitgevoerd en ontmaskerd.}
De vier stukken van het vierkant $8 \times 8$ zijn: twee \omterm{def:g6:shapes:triangles}{rechthoekige driehoeken} met
rechthoekszijden $8$ en $3$, en twee rechthoekige trapezia met parallelle zijden $5$
en $3$ en hoogte $5$.
\begin{enumerate}
  \item Reken de \omterm{def:g6:measure:area}{oppervlakte} van het oorspronkelijke vierkant en die van de beweerde
        rechthoek uit. Formuleer het schandaal in één regel.
  \item Reken de \omterm{def:g6:measure:area}{oppervlakte} van elk van de vier stukken uit
        (\cref{thm:g7:areas:triangle}; voor de trapezia knip je ze door of gebruik
        je \cref{exo:g7:areas:9}).
  \item Tel de vier \omterm{def:g6:measure:area}{oppervlaktes} op. Welke van de twee figuren --- vierkant of
        rechthoek --- kunnen de stukken echt vullen?
  \item In de rechthoekopstelling zou de schuine zijde van een driehoek (die $3$
        stijgt over $8$) de schuine rand van een trapezium (die $2$ stijgt over $5$)
        in één rechte ``diagonaal'' moeten voortzetten. Test die uitlijning: zijn $3$
        over $8$ en $2$ over $5$ dezelfde steilheid
        (\cref{thm:g7:prop:cross})?
  \item Leg uit waar het vijfenzestigste vierkantje zich verstopt: welke vorm heeft
        het gat langs de valse diagonaal, wat is zijn precieze \omterm{def:g6:measure:area}{oppervlakte}, en waarom
        ziet het oog het niet? (Hoe breed is een spleet met die \omterm{def:g6:measure:area}{oppervlakte} gemiddeld,
        uitgerekt over de diagonaal van de rechthoek, die ruwweg $13$ eenheden lang
        is?)
\end{enumerate}

\medskip
\noindent\textbf{Deel II --- Eerlijk redeneren met \omterm{def:g6:measure:area}{oppervlaktes}.}
\begin{enumerate}[resume]
  \item Teken twee \omterm{def:g6:lines:perp}{parallelle lijnen} op $4$ cm van elkaar, een basis van $6$ cm op een
        van de twee, en drie verschillende driehoeken op die basis met hun top op
        verschillende punten van de andere \omterm{def:g6:lines:objects}{lijn}. Reken de drie \omterm{def:g6:measure:area}{oppervlaktes} uit.
        Waarom zijn ze allemaal gelijk, waar de top ook langs zijn \omterm{def:g6:lines:objects}{lijn} schuift?
  \item Bewijs de trapeziumformule in het algemeen: een trapezium met parallelle
        zijden $B$ en $b$ en hoogte $h$, langs een diagonaal geknipt, geeft twee
        driehoeken met dezelfde hoogte $h$. Besluit:
        \[
        A = \frac{(B + b) \times h}{2} .
        \]
  \item Pas haar toe op een trapezium met $B = 9$ cm, $b = 5$ cm en $h = 4$ cm.
  \item Schrijnwerkers gebruiken dezelfde formule vermomd: \omterm{def:g6:measure:area}{oppervlakte} $=$
        (gemiddelde van de twee parallelle zijden) $\times$ hoogte. Leg uit waarom dat
        dezelfde regel is, en reken vraag 8 zo opnieuw uit.
  \item Kantel een nette stapel speelkaarten tot een scheve stapel. Leg uit waarom het
        \omterm{def:g6:measure:volume}{volume} niet is veranderd, en wat dat zegt over de hoogte in de prismaformule
        $V = B \times h$ (\cref{thm:g7:areas:volume}): welke hoogte moet je bij een
        scheve stapel nemen --- de schuine lengte of de verticale? (Vergelijk
        \cref{exo:g7:areas:2}, één dimensie hoger.)
\end{enumerate}

\medskip
\noindent\textbf{Deel III --- Ringen, pizza's en de terugkeer van de goochelaar.}
\begin{enumerate}[resume]
  \item Een ronde vijver met \omterm{def:g3:shapes:shapes}{radius} $3$ m ligt in het midden van een rond grasveld met
        \omterm{def:g3:shapes:shapes}{radius} $5$ m. Reken de \omterm{def:g6:measure:area}{oppervlakte} van de grasring uit (precies met $\pi$, en
        daarna afgerond op de m$^2$).
  \item Zoek de \omterm{def:g3:shapes:shapes}{radius} van de ene schijf waarvan de \omterm{def:g6:measure:area}{oppervlakte} precies gelijk is aan
        die van de ring. (De drie radii die je nu hebt, vormen een beroemd drietal
        waarvan het verhaal in \cref{ch:g8:pythagoras} verteld wordt.)
  \item Wat is meer pizza: één pizza met \omterm{def:g4:geometry:circle}{diameter} $40$ cm, of twee pizza's met
        \omterm{def:g4:geometry:circle}{diameter} $28$ cm? Reken uit en vergelijk.
  \item Met welke factor moet de \omterm{def:g3:shapes:shapes}{radius} van een pizza ruwweg groeien om haar
        \omterm{def:g6:measure:area}{oppervlakte} te \emph{verdubbelen}? Test de kandidaat $1.4$ (denk eraan dat
        \omterm{def:g6:measure:area}{oppervlaktes} het \emph{kwadraat} van de schaalfactor volgen,
        \cref{pb:g6:propdata:1}), en zeg waar de precieze factor --- het getal
        waarvan het kwadraat $2$ is --- zijn officiële entree maakt
        (\cref{ex:g8:pythagoras:diagonal}).
  \item De goochelaar komt terug: hij knipt een vierkant van $13 \times 13$ waarin de
        getallen $5$, $8$ en $13$ de rollen spelen die $3$, $5$ en $8$ eerder speelden,
        en legt er een rechthoek van $8 \times 21$ van. Reken beide \omterm{def:g6:measure:area}{oppervlaktes} uit:
        wat verdwijnt er deze keer? De getallen $3, 5, 8, 13, 21$ zijn opeenvolgende
        termen van de rij uit \cref{pb:g7:fractions:1} --- formuleer het patroon van
        de winsten en verliezen van de goochelaar, en de moraal: waarom liegen
        \omterm{def:g6:measure:area}{oppervlaktes} nooit?
\end{enumerate}
\end{problem}
